\chapter{Bilder und Grafiken}

Bilder bzw. Grafiken lassen sich auf verschiedene Arten mit LaTeX einbinden. 
Es gilt zu beachten, dass im Text ein inhaltlicher Bezug mithilfe von \begin{verbatim}\ref{label_name_grafik}\end{verbatim}
hergestellt ist.\\
Beispiel: In Darstellung \ref{grafik_1} ist das Logo der HHU zu sehen, welches \ldots\\
Viele Dateiformate werden unterstützt, jedoch ist die Nutzung von jpg, png, gif oder pdf (Vektorgrafiken) zu empfehlen.

\begin{figure}
  \centering
  \begin{subfigure}[b]{1.0\textwidth}
    \includegraphics[width=1.0\linewidth]{img/hhu_logo.jpg}
  \end{subfigure}
  \caption{Einzelne Grafik (centered)}
  \label{grafik_1}
\end{figure}

\begin{figure}
  \centering
  \begin{subfigure}[b]{0.5\textwidth}
    \includegraphics[width=1.0\linewidth]{img/hhu_logo.jpg}
  \end{subfigure}
  \caption{Einzelne Grafik (verkleinert)}
  \label{grafik_2}
\end{figure}

\begin{figure}
  \centering
  \begin{subfigure}[b]{0.45\textwidth}
    \includegraphics[width=1.0\linewidth]{img/hhu_logo.jpg}
    \caption{HHU Logo links}
    \label{grafik_3}
  \end{subfigure}
  \begin{subfigure}[b]{0.45\textwidth}
    \includegraphics[width=1.0\linewidth]{img/hhu_logo.jpg}
    \caption{HHU Logo rechts}
    \label{grafik_4}
  \end{subfigure}
  \caption{Zwei Grafiken horizontal angeordnet}
  \label{zwei_grafiken}
\end{figure}

\begin{figure}
  \centering
  \begin{subfigure}[b]{0.8\textwidth}
    \includegraphics[width=1.0\linewidth]{img/hhu_logo.jpg}
    \caption{HHU Logo oben}
    \label{grafik_5}
  \end{subfigure}
  \begin{subfigure}[b]{0.8\textwidth}
    \includegraphics[width=1.0\linewidth]{img/hhu_logo.jpg}
    \caption{HHU Logo unten}
    \label{grafik_6}
  \end{subfigure}
  \caption{Zwei Grafiken vertikal angeordnet}
  \label{zwei_grafiken_2}
\end{figure}
